\begin{definition}[Left-Hand Derivative]
Let $I \subseteq \mathbb{R}$ be an interval, let $f : I \to \mathbb{R}$, and let $c \in I$ be such that $c$ is a left-hand limit point of $I$. We say that $f$ has a \emph{left-hand derivative} at $c$ if the limit
\[
\lim_{x \to c^-} \frac{f(x)-f(c)}{x-c}
\]
exists. In that case, the value of this limit is denoted by $f'_-(c)$, that is,
\[
f'_-(c) := \lim_{x \to c^-} \frac{f(x)-f(c)}{x-c}.
\]
\end{definition}

\begin{definition}[Right-Hand Derivative]
Let $I \subseteq \mathbb{R}$ be an interval, let $f : I \to \mathbb{R}$, and let $c \in I$ be such that $c$ is a right-hand limit point of $I$. We say that $f$ has a \emph{right-hand derivative} at $c$ if the limit
\[
\lim_{x \to c^+} \frac{f(x)-f(c)}{x-c}
\]
exists. In that case, the value of this limit is denoted by $f'_+(c)$, that is,
\[
f'_+(c) := \lim_{x \to c^+} \frac{f(x)-f(c)}{x-c}.
\]
\end{definition}

\begin{theorem}[Differentiability and One-Sided Derivatives]
Let $I \subseteq \mathbb{R}$ be an open interval and let $f : I \to \mathbb{R}$. Then $f$ is differentiable at $c \in I$ if and only if both one-sided derivatives $f'_-(c)$ and $f'_+(c)$ exist and are equal. In that case,
\[
f'(c)=f'_-(c)=f'_+(c).
\]
\end{theorem}






















